\section{重建地方世界：制度在断裂后的落点}
\label{sec:synthesis-rebuilding-local}

帝国重建常被写成一个中心重新占有都城、颁布诏令的时刻。对县邑、关塞和迁徙中的家庭而言，重建从来没有这么整齐。它发生在道路是否再次通行、仓中是否尚有可发之粮、旧吏是否愿意继续办公、邻近的武装是否接受新的名义、逃散的人能否被接纳的许多小决定中。西汉初年的恢复、新朝末年的破坏、东汉的复兴和汉末区域政权的形成，都可以从这些地方性的落点连读。

\subsection{战后第一年：粮食、官印与仍未归来的人}

战争结束并不会立即把人送回田地。楚汉战争之后，一座县城即使重新悬挂汉的官印，仍须面对荒田、残缺的劳力、断裂的亲属关系和不确定的道路。朝廷的任命可以确认谁有权征收和裁断，不能让仓廪自行充满，也不能替已经流散的人决定是否回来。早汉强调休养生息，正因为新政权需要避免把刚刚恢复的家庭再次推入无力承担的征发。这样的选择不仅是仁政语言，也是一种对国家自身税源、兵源和地方服从的判断。

地方官在这段时期的作用不宜只用“执行中央命令”概括。他们需要辨认哪些人仍在本县，哪些土地可以复耕，哪些功臣、宗族和旧有势力能够合作，哪些武装尚未完全解除。史书常保留郡守、相国或诏令的名字，很少连续记录县吏如何逐户处理空缺。正因如此，读者不应以为行政重建只是中央意志从上而下的投影。能否让小规模的交换、婚姻、借贷和互助恢复，往往决定一项大政策在几年后是否仍有可征发的人口。

封国与郡县并存，在这种地方现实中有其功能。它让功臣、宗室和不同区域获得可预期的位置，减少开国分配立即转为全面冲突的可能；它也使资源、武力和人事在部分地区保留了较强的区域性。制度的好处与风险因此同时出现：它为恢复争取时间，却使后来中央必须面对谁有权动用地方兵粮的问题。七国之乱并非从这种安排之外突发，而是将其长期未解决的核验问题推到战场上。

\subsection{新朝的县邑：改革抵达时，谁能解释它}

王莽的改革使地方重建面对另一种困难。西汉末年已有较成熟的官名、货币、市场和地方保护关系；新朝改变名目、币制和土地语言时，县邑中的人必须判断新规定如何与手头的债务、交易和亲属安排相容。诏令可以规定什么叫作合法的钱币、什么叫作新的官职，却不能在同一日替市场建立信任，也不能保证每一位地方官以同样方式解释命令。

这不是把改革的后果完全归为“民众不理解”。不同地区的条件不同：水患、饥荒、道路、地方武装和大姓网络会决定一项新规定在哪里首先遇到阻力。一个依赖货币交易的城镇与一个更依赖实物交换的乡里，面对币制变化的方式不可能相同；一个有强大地方保护者的县与一个刚经历流民涌入的县，对土地和赋役的重新登记也会有不同余地。材料不允许我们逐地复原这些选择，正应使我们避免把“天下大乱”写成一种没有地方差异的反应。

新末的绿林、赤眉及其他地方力量，不只是对首都失败的被动回声。它们在各自可取得粮食、保护同伴和利用道路的条件下组织起来。某些集团能够使用“复汉”或旧有王号的语言，是因为国号、礼制和地方记忆仍有动员价值；可是，名义若不能与城邑、仓储和人群的安全结合，便无法稳定地成为新秩序。新朝的断裂使人看见，制度语言可以快速改变，地方世界的信任与资源却有自己的时间。

\subsection{雒阳复兴：让旧名义重新有用}

东汉初年的复兴并不只是以刘氏名义替换王氏名义。雒阳成为首都后，朝廷需要让地方相信这里发出的任命、赦令和征发不再是转瞬即逝的战时文书。官员能否安全赴任，驿路能否传递奏报，仓储能否支持赈济，都会决定“汉”在一县一乡的实际含义。建武度田正是在这种背景下出现：国家要重新知道谁拥有或耕作土地、谁能够承担赋役，才可能把复兴从政治宣告变成日常治理。

度田也把重建的代价暴露出来。战后土地的占有未必与旧名册一致，逃亡者留下的空缺、地方大户的保护、临时安置的流民，都使“如实登记”成为利益问题。朝廷希望数字清楚，基层却要在相互依赖与相互防备中提供数字。反弹说明登记会触及家庭和地方精英的切身安排；它不能证明所有地方同样拒绝，也不能使朝廷的数字失去意义。两者合在一起，恰好显示国家可见性不是天然存在，而是通过协作、压力与妥协建立。

雒阳的知识网络也是重建的一部分。太学、经师、察举和官员往来，使地方士人能够将名声和学问带到首都，再把首都的评价带回地方。这样的网络为东汉提供人才与共同语言，却使地方声望更容易进入宫廷人事。后来党锢冲突中的名士、学生与官员，并非突然出现的另一种力量；他们正是东汉制度长期鼓励的学习、举荐和评价关系，在宫廷不愿接受某些评价时所显出的政治后果。

\subsection{边地的重建：迁来的人不是空白劳力}

河西、羌地和北方边郡的重建比内地更难以用统一标准衡量。城塞、屯田和传舍需要人，但被迁来的人并不是可以任意安放的劳力。有人带着家庭、技能和旧有关系，有人可能是戍卒、罪人、降人或逃避战乱者；不同身份决定他们能够获得什么土地、承担什么守卫、与当地群体发生怎样的关系。朝廷的“实边”语言表明政策希望达到的结果，不能代替这些人如何适应水源、气候、市场与邻近部众的历史。

西汉的河西经营将军事通道转化为日常行政网络，需要的不是一次胜利而是不断补给。东汉的羌地动员则更明显地表现出供给与地方关系的紧张：当军费和粮运无法稳定，州郡、将领和地方大族必须寻找临时办法。临时办法能让一处关塞不至于立刻失守，却也可能把征发、安置和武装招募的权力集中在区域网络中。边地重建的反馈最终会回到内地：流民进入新县，税粮改道，地方军政人物取得新的谈判地位。

因此，不能把“迁徙”仅作为战争章节末尾的一项结果。它改变家庭与户籍的关系，改变市场与土地的劳力来源，也改变谁可以代表地方与国家交涉。迁徙者不是完全脱离制度的人；他们常在新的地方被重新登记、被安置到屯田、进入依附关系或参加军队。正是这些不同的落点，使同一次战乱在不同地区留下不一样的社会后果。

\subsection{汉末：地方秩序不是中央秩序的简单反面}

黄巾起事、凉州动乱和随后的群雄竞争，使许多县邑不得不在中央有效保护难以抵达时寻找其他安全来源。坞壁、宗族、地方武装、州郡官员和宗教网络可能彼此交错；不能把它们一律视为“割据势力”，也不能把它们都写成保护乡里的善意共同体。它们首先回答的是眼前的问题：谁能守住粮食，谁能在道路不通时传递消息，谁能对外来武装谈判，谁能安排无处可去的人。

地方秩序的出现不意味着皇帝名义毫无作用。曹操奉迎献帝后，诏令、任命和朝廷仪式仍能给北方的行动提供优先的公共语言；这种语言之所以有效，是因为它与屯田、仓储、军队和道路结合。其他区域中心也要使用类似的制度遗产，却会在荆益、江东或关中不同的资源条件中改写它。汉末的断裂不是国家与地方的简单对立，而是不同地方秩序不再愿意或不能够把各自资源送进同一套结算。

\subsection{一册簿籍的两面：地方如何把一家人说成可治理的对象}

地方重建的第一步往往不是修好城墙，而是让某些人的姓名、田地和责任重新能够被写下。县廷的簿籍需要把户、口、田和役连在一起，才可以决定谁应纳租、谁须服役、谁有资格领受赈给；可是一场战乱之后，正是这些关系最不整齐。原来的户主可能已经死亡，成年子弟可能在外从军，寡居者可能与亲属合居，来者可能借住在旧户名下。若把登记当作纯粹技术，便会看不见地方官必须在熟人关系、税役压力和上级期限之间作出的判断。每一次把复杂生活压成一行文书，既使治理成为可能，也留下了有人被漏掉、有人被替代、有人被迫承担他人风险的可能。

西汉开国的恢复尤其显示这种两面性。战后朝廷需要知道关中和各郡还有多少可以耕作、服役和运输的人，也需要避免过早把刚安定下来的家庭再推入逃亡。关于轻徭、减赋、赐爵和迁徙的帝纪记载，能让我们看见政策把人口、土地和财政一起考虑；《汉书·食货志》的后出整理则提供了制度语言，却不等于逐县的执行日志。我们可以说，恢复中的户籍是税源与兵源的基础；不能说某项诏令一出，所有荒地便有了主人，所有户口便准确归册。地方的真实工作，是在播种、借贷、婚姻和邻里互保已经改变的情况下，重新形成一个勉强可用的共同账目。

这个账目也解释了为什么中央会关心迁徙而不只关心征收。前一九六年前后迁关东豪族入关中的主张，把京师安全、人力配置和地方控制放在一起。它并不是把人口从一处搬到另一处即可完成的棋局：迁入者要取得住处、田土与保护，接纳地要重排水源、市场和役使，原籍也会失去或重组原有的中介网络。相关正史记载的十余万口是政策叙事中的重要数字，却有执行范围的争议，不能作为精确人口普查。比数字更可把握的是政策的方向：国家把人户视作可重新配置的资源，而这种配置只有被地方关系吸收时才会成为持续秩序。

\subsection{救荒的路径：从诏令到一碗粮之间的地方权力}

赈济最能检验簿籍是否仍与地方生活相接。诏书可以命令开仓、减租或贷种，却不能自己判断哪一处道路被水冲断，哪一仓的谷物已被军队先取，哪一家虽在旧簿中却已离散。西汉帝纪和《食货志》留下朝廷处置灾害与劝农的记录，主要证明中枢把仓储、赋役和农时当作政治问题；它们很少保存一个县如何排队、由谁验明、是否照顾客居者的细节。史料沉默不应让我们凭想象补出整齐的赈场，也不应使我们误以为灾年只有皇帝和灾民两端。县吏、里中有力者、运输者、仓吏和能够作保的邻人，都可能决定谷物是否越过最后一段路。

对家庭而言，赈给与登记从来互相牵连。领取粮食可能需要证明居处或户名，愿意接受官府安置也可能意味着将来要承担赋役；没有被登记的人则更依赖亲族、市场、雇佣或地方保护。灾后若一地仅按旧簿给粮，流入者容易被排除；若完全不问来历，又可能使有限的仓储和治安承受更大压力。地方并没有一种不会伤人的选择。它只能在谷物数量、播种时日、道路安全和上级考核之间权衡，而这种权衡会改变人们对官府是否值得信赖的判断。

王莽时期特别提醒我们，救荒不能从货币和制度语言中抽离出来。《汉书·王莽传》所叙改币、改制与灾乱，带有后朝叙事的强烈批判，不能直接当作每个乡里的现场纪录；但它足以提示一种实际难题：当交换媒介、官名和产权称谓频繁变动，地方即使有粮也更难以用稳定规则组织借贷、买卖和征发。缺粮者需要的不是抽象的“新制”，而是可辨认的凭据、可通行的道路和不会在次日失效的交换方式。新末的军事集团由此能在若干地区成为给养和保护的替代节点，未必因为人人接受其政治纲领，而是因为它们在国家结算失灵处提供了其他、同样有代价的结算。

东汉复兴后，赦令、复业、度田和赈给应被看成同一组修复工作。赦令使部分人不必因战乱中的旧罪或归附问题永远处于危险；复业与给养让耕作有重新开始的机会；度田则试图把已经变化的人地关系送回赋役的视野。《后汉书·光武帝纪》所见的诏令与度田争议，不能让我们量化所有受惠者，却足以显示地方恢复不是只靠任命官员。若谷物、种子、劳力和文书不能在一个农时内重新接上，恢复的国号就难以在乡里变成可预期的生活。

\subsection{从乡里到州郡：东汉地方中介为何越来越不可替代}

东汉前中期的地方秩序并非中央命令的被动终点。察举、学校、郡国行政和市场道路，使能书写、能担保、能调停的人在县乡与雒阳之间往返。这样的中介对灾后登记十分有用：他们熟悉户口、田界与声望，能帮助官府区分暂居者、依附者和可征发者；但同一份知识也让他们有机会保护自家、隐匿劳力或将公共事务变成私人恩惠。把这一现象简单称为“豪强坐大”，会遗漏地方中介解决的实际问题；把它描述为自然的乡里互助，又会遮蔽被排除的人。重要的是，中介权力随着户籍、赈济与治安责任一同增长。

度田所遇到的阻力，正可放在这个脉络中理解。朝廷要的不是纯粹的数字，而是能据以征收和动员的数字；地方提供的也不是纯粹的信息，而是经过亲属、佃作、寄居和保护关系过滤的信息。若一户被多报或漏报，后果会落在赋役、赈贷和邻里责任上，因此任何登记都可能成为冲突的焦点。现代研究常以人口数字衡量国家能力，但《后汉书》的帝纪、志书和列传并不是同一时点、同一口径的统计表。它们更适合用来追问：国家在哪些地方还能够要求地方说明人户，地方又在哪些条件下能够把自己的说明变成事实。

边郡使这一问题更为尖锐。河西、北边与羌地的城塞需要运输、守卫和屯垦，迁入者与原住者却未必愿意承担相同风险。简牍能让我们看见某些关塞的出入、供给与文书程序，说明帝国行政确实抵达过具体站点；它们不能代表所有边地家庭，更不能把军政安置写成没有冲突的成功。地方官和军政人物在筹粮、安置、雇募时愈发依赖本地网络，网络便取得以“保境”“济急”名义决定资源流向的能力。这种能力在平时可补足中枢的距离，在危机中也可能使州郡和强宗优先保护自己的结点。

\subsection{汉末的分裂账本：谁还能决定一家人的归属}

黄巾之后的连续动乱，使县廷原有的簿籍、仓储和治安逐步失去同一节奏。并非每个地方都在同一时间完全失序：有的县仍可依据官府名义征发，有的地方靠坞壁、宗族或州郡军队维持粮食，有的道路则先被商旅、逃亡者和兵队重新定义。史书多用“流民”“盗贼”“部曲”等概称，便于叙述军政变化，却会掩盖同一称谓下不同的家庭处境。逃离兵火的人可能在亲属处暂住，可能被招为兵，可能进入大族庄园，也可能被屯田或区域政权重新编制；这些落点决定了他后来缴粮、服役和寻求保护的对象。

坞壁与宗族不是中央秩序的简单反面。它们能在道路不通时储粮、守望、向军队谈判，也能把外来者的劳力纳入依附和债务。对寄居家庭来说，得到一处围墙内的安全，常以接受新的担保和约束为条件；对拥有仓谷和武装的人来说，接纳人口既能增强生产与防卫，也会增加分粮和报复的风险。因而所谓地方碎裂，不是地方社会突然摆脱国家，而是原先由县廷统一协调的一部分登记、赈济与保护，分别被不同的区域节点承担。它们仍使用官名、印信和文书，却不再向同一个中心交付相同的优先次序。

曹操的屯田可见这种重新结算的努力。《三国志·魏书·武帝纪》与《后汉书》末年材料能证明北方政权把流民、荒地和军粮放在紧密关系中；后世概括的成效不能代替各地安置的细部。屯田给某些家庭提供耕作和口粮的路径，同时把生产、迁徙和军事纪律捆得更紧。它不是对家庭危机的单纯救济，也不是只为战争的冷酷技术，而是一种以新的户籍和劳动约束换取可持续供给的安排。区域政权能否在这种安排中守住道路与仓储，决定它能否把暂时的收容转为长期的税源和兵源。

当献帝禅位时，旧的县邑和文书传统仍在，消失的是一个能够让各处簿籍、粮仓与征发最终对同一朝廷负责的稳定中心。魏、蜀、吴各自要面对流民、寄居、屯垦和地方保护，却会以不同的军政与资源条件作答。将汉末理解为“地方化”，不是给它贴上衰败标签，而是看到国家原来在家庭尺度上提供的几项功能——登记、赈给、裁断与安全——已被分散的权力分别接管。家庭不是宏大崩解的旁观者；他们每一次借粮、投靠、迁出或重新入籍，都是新地方秩序得以成立的材料。
\subsection{重建的制度得失}

每一次重建都解决了真实的难题：早汉以多种地方安排恢复战后协作，新朝试图重订已经失衡的分类，东汉重新使雒阳的文书和户籍发生作用，汉末各地则在危机中维持最低限度的供给与安全。代价也同样真实。恢复越依赖地方中介，地方中介越可能将公共资源变成私人保护；登记越追求清晰，越可能将流动、依附和家庭内部差异压成不合实际的数字；边地越需要临时军政网络，越难保证网络在危机后仍只服务中枢。

将这些得失放在同一条跨期线索中，并不是要找出一个导致汉亡的单一原因。它让读者看见，西汉、新、东汉不断处理的是相似而不相同的连接问题。汉帝国的延续来自能够重新接线的能力；它的终结则发生在重接的代价不再能由一个中心共同承担时。

\subsection{三次开仓：同一动作并不产生同一种恢复}

汉代的诏令常以开仓、赈贷、蠲租或遣使巡行来显示朝廷看见了灾年。把这些动作
排成一条从西汉到东汉的进步或衰败曲线，都太过轻易。西汉初年面对的是战后田土、
人口和道路尚未重新接合的局面；新莽末年则是币制、官名和政治归属本身都变得
难以取信；东汉的赈给发生在重建户籍与地方中介的同时，到了汉末，粮食又常先被
卷入守城、募兵和流民安置。仓门可以在不同年代被命令开启，仓外等待它的人、
负责核验它的人和能够将粮带走的道路却已经不同。

西汉有关劝农、赈贷和减免的记载，足以说明新政权把仓储与农时视为恢复人户的
条件。它们并没有保留每一县如何称量谷物、谁先取得种子、谁因缺少担保而被排除。
这一限度尤其重要，因为早汉的“休养生息”很容易被写成一种从长安平滑施到乡里的
政策。实际上，地方要先有尚能运作的里、县、市场和人际关系，较轻的赋役才可能
变成一季真实的喘息。若战后家庭已经失去劳力或牲畜，减轻征收并不会自动修复
耕作；它只是使重新开始不至于立刻被下一次征发打断。

新朝的困难又不在于完全没有救荒语言。王莽曾遣人开仓、赈贷穷乏，传世叙事还
记有教民煮木为酪等处置。后者带着强烈的失败叙事，不能让我们替所有地方复原
一种荒诞场景；可是它确实指出，当谷物、货币、道路和官吏之间的正常交割已被
破坏，发出救济命令本身并不足够。若地方不能相信新币可以换物、令文可以持续、
仓粮不会被中途挪用，那么救助会加重而非解除不确定性。新莽的断裂提醒读者，
制度语言越想迅速覆盖全国，越依赖基层已有的信用与物资。

东汉早期的复业、赦令、度田与赈给，则把恢复放到更明确的簿籍工作中。国家需要
知道谁还在、谁耕作何地，才谈得上将谷物、赋役与兵役重新分配；而战乱后的家庭
又恰恰最难以旧日身份回答这些问题。地方大户、县吏、亲属和行旅者于是成为不可
或缺的中介。中介可以帮助国家发现灾情和安置来者，也可能选择性地保护自己的人。
这不是东汉独有的道德缺陷，而是所有大规模救助面对的结构性难题：救助需要地方
知识，地方知识又不能轻易被中央完全接管。

\subsection{道路被修好以后，归属并不会自动归位}

道路在恢复叙事中常被当成中性的背景。其实，一段道路重新可走，首先改变的是谁
能够带着粮食、文书、兵器、家人或债务穿过它。楚汉战争后的关中道路让新朝能够
任命官吏、征集资源，也让迁入者和商人重新调整去处；河西经营把通向西域的路线
接入帝国，同时使泉水、传舍和沿线家庭承担不断的供给；新末、东汉初年和汉末的
道路破坏，则让同一片地方的安全、价格和信息各自缩回更近的节点。道路不是一条
天然属于皇帝的线，它必须由守卫、运输、修缮和人们愿意使用的凭验组成。

因此，迁徙也不能只放在战争的结尾。早汉迁豪族入关中、边地安置戍卒与降附者、
新末流民和东汉羌乱中的迁离，分别由不同的政治目的和生计压力推动；它们共同
要求接纳地处理土地、水源、登记和治安。被迁者并非一团可任意填进空地的劳力，
接纳者也并非自然愿意分享资源。国家有时用赐爵、给田、免役或军政编制使安置
显得可行，实际的融入仍取决于地方关系是否承受得住。材料少有连续记下一个家庭
从出发到安顿的全过程，这不是把他们消失在“徙民”二字中的理由。

边地最能显示道路与归属的纠缠。西汉关塞文书保存传符、给马和出入的细碎动作，
东汉正史则更多从都护、将领与战事讲述西域、凉州和北边。两种材料不能拼成一份
连续的帝国总账，却能互相限制我们的想象：一方面，行政确实要靠核验、供给和
人员交接；另一方面，不能把西汉某一关塞的程序直接复制为所有东汉边郡的日常。
所谓实边、屯田和守塞，只有在具体水源、农时、市场和邻近关系中才有意义。

\subsection{汉末的“地方”，是许多被迫接手的接口}

汉末地方化并不意味着中央秩序突然消失，或每一处乡里都自愿转向独立。献帝的
诏令、官名和礼仪依然可以被不同集团使用，因为它们仍是安排任命、结盟和公开
承诺的有效语言。变化在于，谁能让这些语言同粮仓、守卫、文书与人口重新接上，
谁才有能力使它们在一地生效。州牧、郡守、坞壁主、宗族领袖、将领和宗教网络
之所以在史书中愈发显眼，不只因为他们更有野心，也因为他们被要求接手原先需要
跨地区协调的事务。

黄巾与凉州的危机把这一过程推得格外清楚。县城需要守卫，军队需要粮食，逃亡者
需要住处，地方必须在消息混乱中判断谁能作保。一个富室出粮、一名将领临时募兵、
一处坞壁收留流民，都可能暂时补上国家来不及处理的缺口；同时，这些行为也把
债务、人身依附、军功和任命集中到新的节点。它们不是皇帝名义的纯粹反面，往往
还要借诏令和官印取得正当性。真正分散的是结算的终点：不同地区不再愿意或不能
把粮、人和武力的优先次序交给同一个中心。

曹操的屯田、江东的水路军政和益州、荆州等地的地方安排，都应放在这个问题下
理解。它们不是凭空出现的“军阀制度”，而是对流民、荒地、运输和保护作出的
不同接法。每一种接法都可能给部分人提供相对稳定的耕作与安全，也会让另一些人
接受新的役使和归属。到二二〇年，禅代改变了最高名号，却没有把这些接口还原为
旧汉的统一网络。魏、蜀、吴继承的不是一片等待重新命名的空白，而是已经由战争、
救助、迁徙和地方协商重新编织过的社会。

\subsection{文书重新抵达之前：恢复总是有不同的速度}

一座县城在战后重新收到印信和文书，并不表示它已经重新拥有帝国。印信能够让
官吏以某个名义征收、裁断和发给凭验；它不能在一日之内带回逃散的人口、修好被
冲坏的桥梁，或使隔了数年的债务重新有了双方都接受的尺度。西汉开国、东汉复兴
和汉末各区域政权建立，都有“文书重新抵达”的时刻。它们的共同点不是新政权
突然拥有无所不至的力量，而是先把若干已经断开的交接——官吏与地方、仓储与
农时、道路与消息、户名与家庭——尝试接回去。

这种接回有先后。都城或州治往往先能恢复任命和传令，因为少数人、印章和文书
即可构成一个可见的中心；较晚恢复的，是要许多家庭同时参与的耕作、市场和互信。
一条命令从雒阳或长安传到县廷，可以要求重新登记和征收，县廷却不能凭命令得知
一名外出从军的儿子是否回来、一块田是否已由邻人代耕、一个寄居者是否愿意承担
新的赋役。行政最初得到的往往是旧名册、临时的保证和相互矛盾的口供。国家不是
从一张空白纸开始，而是在已经变化的生活上重写可用的分类。

这也使“恢复户籍”成为一件需要谨慎理解的事。簿籍既是国家重新看见税源、兵源
与赈给对象的工具，也是家庭争取归属、避免被误征或取得救助的入口。一家人被写
回某个县，并不必然意味着他们回到了战前的土地和关系；有时只是承认他们已经在
别处活下来。另一方面，未被写入的人并非不存在，只是更难通过官府的语言主张
自己。传世史料常以户口数字、诏令和地方官名显示重建已经发生，较少留下这些
边缘身份如何被协商。正因如此，数字可以说明国家想要达到的可见性，不能替代
一幅完整的灾后社会照片。

王莽末年的经验将这种速度差推到极端。旧汉的官名、钱币和地方保护关系尚未消失，
新朝的称谓、币制和行政命令又要求人们按另一种尺度交割。后朝史家把新莽写成
覆亡的前史，材料的立场要求我们不能把每个县的混乱预先写成结局；但制度更新
必须经过日常交接这一点无可回避。一枚钱能否被收下，一张文书能否被相信，一个
地方吏能否调动仓粮，取决于人们是否预期下一次交割仍有同样的规则。当这个预期
消散，最先填补空缺的往往不是宏大的反叛宣言，而是亲属、市场、武装和地方名望
提供的较近而较窄的承诺。

\subsection{收容不是终点：流民抵达以后发生的政治}

“流民入某郡”在编年材料里通常只占一行，真正的难题却从抵达以后开始。接纳地
要面对的不是一群抽象的人口，而是带着亲属、疾病、债务、耕作经验和旧有冲突的
家庭；原有居民则要衡量水源、田地、劳力和治安会怎样改变。官府可以设法给粮、
分配荒地或暂时免役，地方有力者也可能提供住处、借贷或工作。每一种收容都使
人得到某种保护，也把他们放进新的责任链：领粮可能需要作保，借种可能形成债务，
得到围墙内的安全可能意味着为守卫或耕作出力。

这不是要把地方收容写成冷酷的计算。战后若没有亲属、宗族、邻里、寺祠、商人或
地方豪右能够先行提供一顿饭、一间屋、一条消息，许多家庭根本等不到官府的安排。
问题是这些帮助无法从权力关系中抽离。能够收留人的一方，会取得谈判土地、劳力
和安全的更大位置；被收留者也会在新的保护者之间选择，或在无法选择时被编入。
国家希望借助这些网络恢复生产和秩序，却又担心网络把人口和资源固定在私人的
名下。两汉反复出现的清查、赦令、迁徙和限制，并不说明国家已经找到了无代价的
答案，恰恰说明它始终在处理这一矛盾。

边地和内地的收容也不应混为一种经验。河西、陇右和北边的来者，可能同时面对
守备、屯戍、牧地与异地市场；关中、山东或荆州的流民，则可能在更密的人口和
地方精英网络中寻找落点。国家文书会以徙民、降人、流民等名词将他们归类，名词
有组织行政的必要，却不能消去差别。有人被安排到军屯，有人进入佃作或部曲，
有人借亲属重入旧籍，也有人在下一次战事中再次离开。把所有移动都写成中央的
人口工程，或都写成个人自由逃亡，都会忽略这些约束。

汉末的坞壁特别容易被浪漫化为地方社会自然的自救，或反过来一概视为割据的
胚胎。较贴近史实的问题是：围墙、粮仓和武装究竟替谁解决了什么。它们能够在
道路不通时保存粮食、让附近人群得到暂时庇护，也能把进入者置于拥有者规定的
劳动、债务和服从之下。县廷的作用若已削弱，坞壁便可能承担登记、分粮和守卫的
一部分；这不表示它已成为一个脱离国家的完整政权。它仍可能借用官名、印信和
朝廷任命来同外部交涉，只是这些名义不再保证资源最终流向同一中心。

\subsection{地方碎裂不是没有秩序，而是秩序不再共用同一结算}

把汉末称为“崩解”，容易让人想象一切规则都消失了。事实上，许多规则仍在被
使用：人们仍需凭文书通行，仍会为土地、婚姻和债务寻找裁断，仍要用某种方式
储粮、征发和安置逃亡者。改变的是谁有能力给这些规则最后的效力。过去一份县
廷文书可以在更广范围内预期得到支持；到汉末，同一张文书的分量要看它背后有
没有兵、粮、道路和可被承认的保护网络。地方不是没有秩序，而是不同地方开始
以不同的节点结算承诺。

这种碎裂也解释了皇帝名义的持久。皇帝并未因为雒阳失守或献帝迁徙而突然失去
一切意义；诏令和任命仍能为竞争者提供共同可读的语言。可是名义无法自行穿过
被截断的道路，也无法替一处仓库补足谷物。董卓、袁绍、曹操以及江东、荆益的
行动者之所以都必须处理地方关系，不是因为他们不懂得正统，而是因为正统若不能
与具体的供给和守卫结合，便无法使陌生人承担代价。皇帝成为资源，正是由于他的
名义仍可组织行动，而组织行动所需的物质条件已落在不同集团手中。

从这个角度回看两汉，地方世界从来不是中央命令之外的静止背景。它保存粮食，
传递消息，安置迁徙者，解释疾病和灾异，也在危机中决定什么样的命令还值得执行。
西汉能够在战后重新整合，东汉能够在新末之后复兴，靠的不是消除地方中介，而是
暂时使中介、道路和簿籍又朝向同一套结算。汉末的终局也不是这些地方联系忽然
出现，而是它们在持续战乱和供给压力下不再能够同时接受一个中心的优先次序。
这条线索不能给汉亡归结为单一制度缺陷，却能让读者看见帝国终结如何发生在仓门、
渡口、户名与保护关系这些最具体的地方。
